\documentclass[preprint,12pt]{elsarticle}
\usepackage{amsmath,amssymb}
\usepackage{graphicx}
\usepackage{booktabs}
\usepackage{hyperref}

\journal{Journal of Energy Storage}

% AUTO-GENERATED numeric macros -- every number cited in this manuscript MUST come from here.
% Generated by make_results_registry.py from results_registry.json. Never hand-type a number
% below this line; if a macro is missing or reads UNVERIFIED, the corresponding claim has not
% been verified yet and must not be stated as fact.
% AUTO-GENERATED by make_results_registry.py -- do not hand-edit.
% Every number in main.tex must come from a \NUM{...} macro defined here.

\newcommand{\MainModelMAE}{0.02479}
\newcommand{\MainModelRMSE}{0.10174}
\newcommand{\MainModelCoverage}{90.5}
\newcommand{\CoverageBIT}{88.0}
\newcommand{\CoverageCALCE}{75.7}
\newcommand{\CoverageMIT}{92.3}
\newcommand{\CoverageNASA}{66.7}
\newcommand{\CoverageOxford}{100.0}
\newcommand{\CoverageRWTH}{90.7}
\newcommand{\CoverageXJTU}{93.0}
\newcommand{\MultiseedMAEMean}{0.02548}
\newcommand{\MultiseedMAEStd}{0.00055}
\newcommand{\MultiseedCoverageMean}{90.33}
\newcommand{\MultiseedCoverageStd}{0.40}
\newcommand{\LodoNASAMAE}{0.00668}
\newcommand{\LodoNASACoverage}{65.5}
\newcommand{\LodoCALCEMAE}{0.01447}
\newcommand{\LodoCALCECoverage}{60.2}
\newcommand{\LodoOxfordMAE}{UNVERIFIED}
\newcommand{\LodoOxfordCoverage}{UNVERIFIED}
\newcommand{\LodoMITMAE}{UNVERIFIED}
\newcommand{\LodoMITCoverage}{UNVERIFIED}
\newcommand{\LodoBITMAE}{UNVERIFIED}
\newcommand{\LodoBITCoverage}{UNVERIFIED}
\newcommand{\LodoXJTUMAE}{UNVERIFIED}
\newcommand{\LodoXJTUCoverage}{UNVERIFIED}
\newcommand{\LodoRWTHMAE}{UNVERIFIED}
\newcommand{\LodoRWTHCoverage}{UNVERIFIED}


\begin{document}

\begin{frontmatter}

\title{Dataset-Aware Conformal Calibration for Cross-Dataset Lithium-Ion Battery State-of-Health Prediction}

\author{Raheeb Sadman}
\address{BRAC University}

\begin{abstract}
% TODO (post-K6): finalize once the repeated-split (K6) between-split variability numbers are
% available -- this draft states the claim at the level of generality the CURRENT verified
% evidence supports, and will be tightened (or loosened) once K6 completes. Do not add a specific
% between-split coverage range here until \RepeatedSplit* macros exist in numbers.tex.
We predict the state of health (SOH) of lithium-ion cells from short windows of historical
capacity-fade data using a multi-scale bidirectional LSTM with self-attention, trained and
evaluated across seven public datasets spanning different institutions, cell chemistries, and
testing protocols (\MainModelCoverage\% pooled groupwise conformal coverage at a nominal 90\%
target, MAE \MainModelMAE, RMSE \MainModelRMSE, $N=339$ cells, 337 contributing
windowed samples). Unlike prior work that either validates distribution-free uncertainty
quantification on a single data source or fuses multiple heterogeneous sources without any
coverage guarantee, we apply groupwise (Mondrian) split-conformal calibration, an empirical-Bayes
partial-pooling (hierarchical) calibration, and a maximum-mean-discrepancy-weighted heterogeneity-
guided conformal predictor (HGC-CP) across all seven sources simultaneously, and report every
quantity from machine-generated artifacts rather than hand-typed values. We find that
calibration-level generalization (held-out calibration data, model trained on all sources) differs
sharply from model-level generalization (the entire source dataset withheld from training,
leave-one-dataset-out): coverage for genuinely unseen sources ranges from
\LodoBITCoverage\% (BIT) to \LodoRWTHCoverage\% (RWTH), well below what calibration-level
partial pooling achieves on the same splits. We additionally show, via independent re-verification
against the raw source files, that two of the seven datasets (XJTU, RWTH) carry previously
undocumented or mischaracterized data-quality artifacts that materially affect a non-trivial
subset of predictions, and we quantify rather than dismiss their effect. We report where our
model fails to outperform a trivial persistence baseline, where a heterogeneity-aware conformal
predictor's hyperparameters cannot be chosen via a genuine held-out protocol without the method
performing poorly on at least one held-out source, and we narrow every claim to what the repeated,
cell-level-valid evidence actually supports.
\end{abstract}

\begin{highlights}
\item Multi-scale BiLSTM + attention model validated across 7 public datasets (339 cells, 337
contributing windows), not a single data source.
\item Groupwise, pooled, normalized-groupwise, hierarchical (empirical-Bayes), and MMD-weighted
heterogeneity-guided (HGC-CP) conformal calibration are compared under a predictor-consistent
protocol, with every reported number traced to a machine-generated registry.
\item Independent re-verification of two datasets' (XJTU, RWTH) documented data-quality
corrections against the current raw files, with a quantified effect of previously-mischaracterized
artifacts (XJTU sentinel-value bridging windows: \XjtuBridgingFraction\%, \XjtuBridgingWindowCount{}
of the XJTU test set, see Section~\ref{sec:data-quality}).
\item Leave-one-dataset-out (model-level) generalization is shown to be substantially harder than
calibration-level partial pooling, and is reported honestly rather than conflated with it.
\item A nested, nothing-tuned-on-test held-out hyperparameter search shows that MMD-weighted
heterogeneity-guided calibration cannot be configured with a single setting that generalizes
coverage uniformly across all seven sources -- reported as an exploratory, not held-out-validated,
result.
\end{highlights}

\begin{keyword}
lithium-ion battery \sep state of health \sep conformal prediction \sep uncertainty quantification
\sep cross-dataset generalization \sep hierarchical calibration \sep data quality
\end{keyword}

\end{frontmatter}

\section{Introduction}
\label{sec:intro}

Accurate state-of-health (SOH) estimation underlies safe and economical operation of lithium-ion
battery systems in electric vehicles and stationary storage, and the great majority of recent
data-driven SOH work reports point-prediction accuracy (MAE/RMSE) without a calibrated,
distribution-free statement of how much to trust an individual prediction
\citep{natcomm2023soh,transformertransfer2025,fastcharge2025}. Where uncertainty quantification is
reported, it is almost always validated on a single public dataset or a single chemistry
\citep{natcomm2023soh}, leaving open whether a calibrated interval's nominal coverage survives
contact with a genuinely different data source -- a different institution's test protocol,
a different cell chemistry, a different cycling regime.

A parallel line of work fuses multiple heterogeneous battery datasets to improve point-prediction
accuracy and generalization
\citep{learningtofuse2025,multitasksohrul2026,batterylife2025}, but none of these report a
distribution-free coverage guarantee for the resulting predictions: multi-source battery fusion
and calibrated uncertainty quantification have, to date, been pursued as separate research
threads. The one paper we are aware of that combines maximum-mean-discrepancy-based domain
adaptation with conformal prediction for battery SOH,
\citet{filgueira2026conformalized} (submitted March 2026), does so on a \emph{synthetic} virtual
battery dataset constructed to simulate manufacturing and usage variability in a controlled way --
an important and complementary contribution, but one that leaves open whether the same approach
survives the uncontrolled, confounded heterogeneity of real public datasets collected by different
labs under different protocols.

This paper closes that gap directly: we apply split-conformal prediction (pooled, groupwise/
Mondrian, and normalized-groupwise), an empirical-Bayes hierarchical (partial-pooling) calibration,
and a maximum-mean-discrepancy-weighted heterogeneity-guided conformal predictor (HGC-CP, in the
sense of \citet{tibshirani2019covariate}'s corrected weighted-conformal quantile and
\citet{barber2023beyond}'s validity analysis for approximate, non-likelihood-ratio weights) to a
multi-scale BiLSTM-with-attention SOH forecaster trained across seven REAL public datasets (MIT/
Severson, RWTH Aachen, XJTU, BIT, CALCE CS2/CX2, Oxford, NASA; 339 cells). We treat every reported
number as something that must be machine-generated from saved model outputs, not hand-typed, and
we report honestly where our own results fall short of what a reader might hope for.

\subsection{Research gaps addressed}

\begin{enumerate}
  \item \textbf{Single-source validation of calibrated uncertainty.} Battery SOH uncertainty
  quantification work to date validates on one dataset or chemistry at a time
  \citep{natcomm2023soh,transformertransfer2025}; we validate groupwise, pooled, hierarchical, and
  HGC-CP calibration simultaneously across seven heterogeneous real sources, and report coverage
  per source rather than only pooled.
  \item \textbf{Multi-source fusion without a coverage guarantee.} Recent multi-source battery
  fusion work \citep{learningtofuse2025,multitasksohrul2026,batterylife2025} improves point
  accuracy or proposes benchmarks but reports no distribution-free interval coverage; we report
  calibrated coverage for every pooling strategy compared.
  \item \textbf{Synthetic vs. real heterogeneity in MMD-weighted conformal calibration.}
  \citet{filgueira2026conformalized} is, to our knowledge, the only other MMD-weighted conformal
  approach for battery SOH, validated on synthetic manufacturing/usage variability; we apply the
  same class of method to real, uncontrolled, cross-institution heterogeneity and report where it
  succeeds and where (Section~\ref{sec:hgccp-results}) a single hyperparameter setting cannot be
  chosen that generalizes coverage uniformly across sources.
  \item \textbf{Model-level vs. calibration-level generalization, routinely conflated.} We show
  these are empirically very different regimes (Section~\ref{sec:lodo}): a model trained on six
  datasets and evaluated, with recalibration, on held-out data from those same six datasets
  behaves very differently from a model that has never seen the seventh dataset's distribution at
  all during training. We report both and never describe the latter as reflecting the former.
  \item \textbf{Window-level vs. cell-level statistical validity.} Successive windows drawn from
  one cell's degradation curve are correlated, not independent trials -- a concern raised at the
  level of general conformal theory \citep{barber2023beyond} but, to our knowledge, rarely
  operationalized with an explicit cell-clustered bootstrap in applied battery SOH papers. We
  report both window-level and cell-clustered-bootstrap confidence intervals throughout, and show
  concretely (Section~\ref{sec:coverage-ci}) that for some sources the two differ by several-fold.
  \item \textbf{Public battery dataset quality artifacts, under-documented and under-tested.}
  Recent reviews confirm that quality and consistency vary substantially across the community's
  public aging datasets \citep{agingdatasetsreview2024}. We independently re-verify, against the
  current raw files, two previously-documented or newly-identified data-quality corrections (XJTU
  sentinel-value bridging, RWTH session-boundary truncation) and quantify their effect on reported
  metrics with an explicit with/without sensitivity analysis, rather than asserting a correction's
  adequacy without re-checking it.
\end{enumerate}

\section{Related work}
\label{sec:related}

\subsection{Deep learning architectures for SOH/RUL estimation}
Recurrent and attention-based architectures remain the dominant approach for capacity-fade
forecasting from cycling data. CNN-LSTM-attention hybrids and multi-scale designs that combine a
raw-resolution branch with a coarser, pooled branch are an active direction for capturing both
short-term noise and longer-term degradation trend. Our model variant ablations
(Section~\ref{sec:ablations}) isolate the contribution of the self-attention pooling layer and the
multi-scale branch directly, rather than only reporting the full model's accuracy.

\subsection{Conformal prediction and distribution-free uncertainty quantification}
Split conformal prediction \citep{angelopoulos2023gentle} gives a marginal, distribution-free
coverage guarantee calibrated on a held-out set, at the cost of a single pooled interval width that
can be miscalibrated for any individual subgroup whose error distribution differs from the pooled
population -- exactly the regime a multi-source, multi-institution battery study produces. Weighted
conformal prediction under covariate shift \citep{tibshirani2019covariate} and its extensions to
approximate, non-likelihood-ratio weights \citep{barber2023beyond} and to posterior drift
\citep{posteriordrift2025,jonkers2024cps} provide the theoretical basis for the
heterogeneity-weighted calibration we apply in Section~\ref{sec:hgccp}, while making clear -- a
point we carry through explicitly -- that a heuristic (non-likelihood-ratio) weight choice does not
inherit an exact finite-sample guarantee merely by using the corrected weighted-quantile formula.

\subsection{Cross-dataset and multi-source battery generalization}
Transfer-learning approaches to battery SOH typically adapt a model from one source dataset to one
target dataset \citep{transformertransfer2025,fastcharge2025}, or select among candidate source
domains by a distributional-distance criterion such as maximum mean discrepancy. Multi-source
fusion approaches combine several datasets' training signal, with entropy- or distance-based
weighting across sources \citep{learningtofuse2025}, or propose a shared multi-task objective
across several datasets at once \citep{multitasksohrul2026}; a recent large benchmark effort
\citep{batterylife2025} assembles a substantially larger pooled dataset across several chemistries.
None of these report a calibrated, distribution-free coverage guarantee for the resulting
predictions. The only directly comparable MMD-weighted-conformal approach we are aware of,
\citet{filgueira2026conformalized}, targets manufacturing/usage variability via a synthetic
dataset rather than real heterogeneous public sources, as discussed above.

\subsection{Hierarchical and partial-pooling calibration across heterogeneous sources}
Outside the battery literature, hierarchical (partial-pooling) models are an established way to
share statistical strength across related-but-distinct groups while still allowing each group its
own estimate, with the degree of pooling determined by the data rather than fixed in advance --
for example in knowledge transfer across structurally similar but distinct engineering fleets
\citep{bull2023hierarchical}. We apply the same logic (an empirical-Bayes, DerSimonian-Laird
random-effects estimator) to conformal calibration quantiles across our seven battery datasets,
rather than to raw point predictions, and validate our implementation against hand-calculated
cases and an independent from-scratch reimplementation (Section~\ref{sec:hierarchical}).

\subsection{Baselines and the persistence baseline specifically}
A trivial one-step persistence forecast (predict the next value equals the last observed value) is
an easy baseline to under-report in applied time-series forecasting, and a recent methodological
study emphasizes that naive forecasts are frequently stronger than they are given credit for in
time series with high persistence \citep{beck2025naiveforecast}. We follow this recommendation
directly: Section~\ref{sec:persistence} reports where our trained model does, and does not, beat a
persistence baseline, including cases where it does not.

\section{Datasets and preprocessing}
\label{sec:data-quality}

We use seven public datasets: MIT/Stanford/Toyota Research Institute's fast-charging cell set
(``MIT''), RWTH Aachen's commercial-cell aging set (``RWTH''), XJTU's battery degradation set
(``XJTU''), a Beijing Institute of Technology fast-charging set (``BIT''), CALCE's CS2/CX2
Arbin-cycled cells (``CALCE''), the Oxford Battery Degradation Dataset (``Oxford''), and NASA's
PCoE battery aging set (``NASA''). SOH is computed as capacity at a given cycle divided by the
median of that cell's first three cycles' capacity (not capacity at cycle 1 alone, which is
noisier). Cells are split train/validation/test (70/15/15) at the \emph{cell} level, never the
window level, so no cell contributes windows to more than one split. Across all seven datasets
this yields \TotalCellsLoaded{} loaded cells, of which \CellsWithWindows{} contribute at least one
windowed sample (one NASA cell and one RWTH cell have too few cycles for even one window at the
input length used); \TrainWindows{} / \ValWindows{} / \TotalTestWindows{} windows in train/
validation/test respectively.

\subsection{Independent re-verification of data-quality corrections}
Three of the seven datasets required non-trivial data-quality correction before use, and we
re-verified each correction against the \emph{current} raw files during this work (not merely
inherited the correction's original justification):

\textbf{CALCE.} Four defects were found and corrected: (i) 331 of 366 raw Excel files record
\texttt{Discharge\_Capacity} as a cumulative counter within the file rather than a per-cycle value,
requiring differencing; (ii) short cumulative-counter files can tie at their last row, requiring an
explicit tie-breaking rule; (iii) one file
(\texttt{CX2\_16\_..\_self discharge test.xlsx}) is a non-cycling self-discharge test and is
excluded; (iv) a final-cycle truncation flag is computed \emph{after} differencing, since the
truncation signature only becomes visible post-differencing. After correction, the loader produces
14{,}188 rows across 10 cells, a per-cell SOH-filter pass rate of 99.46\%, and a median absolute
first-difference of 0.00168 (within the 0.001--0.003 range expected for genuine degradation rather
than noise or alternation). An \texttt{INCOMPLETE\_POLICY=drop} sensitivity check (dropping rather
than merely flagging incompletely-captured cycles) removes 0.63\% of rows; central CALCE
conclusions are unaffected by this choice.

\textbf{XJTU.} Eight \texttt{Sim\_satellite\_*} cells required removal of periodic sentinel/
calibration readings that are not real degradation measurements (detected by frequency: a
handful of distinct capacity values recur far more often than any real degradation trajectory
would produce). We found, independently of this correction's original design intent, that the
surviving rows' cycle numbers are \emph{not} renumbered after the sentinel rows are removed --
only the row index is reset. Consequently every one of the \XjtuBridgingWindowCount{} test windows
drawn from these 8 cells (\XjtuBridgingFraction\% of XJTU's test set) spans a point where sentinel
rows were removed, joining the capacity reading just before a removed block to the one just after
it as though they were consecutive cycles. We quantify rather than dismiss this: these windows show
substantially elevated point-prediction error relative to XJTU's other test windows (see
Section~\ref{sec:xjtu-sensitivity}), while split-conformal coverage validity is not broken by their
presence.

\textbf{RWTH.} Two corrections were independently re-verified against the raw session CSVs:
(i) a session-boundary structural correction -- the raw \texttt{AhEla} (cumulative discharged Ah)
resets at the start of each new session log file, not at the true start of whichever cycle happens
to be in progress, so the first cycle recorded in a session systematically under-reports capacity;
dropping each session's first recorded cycle removes 734 of 114{,}716 rows (0.64\%); (ii) a
statistical local-outlier safety net for isolated single-row artifacts not explained by the
structural correction (e.g.\ a multi-row near-zero stretch in one cell's trace), removing a further
787 rows (0.69\%) when applied alone, or 848 rows (0.74\%) combined with (i) -- less than the sum of
the two because the corrections overlap on a subset of rows.

\section{Methods}
\label{sec:methods}

\subsection{Architecture}
The forecaster is a multi-scale bidirectional LSTM with self-attention: a raw-resolution branch and
a coarser, pooled-resolution branch each process the same input window, are concatenated, and fed
through a self-attention pooling layer before a regression head. Ablation variants remove the
coarse branch (\texttt{no\_multiscale}), replace attention pooling with mean pooling
(\texttt{no\_attention}), collapse to a single plain BiLSTM (\texttt{plain\_lstm}), or substitute a
GRU (\texttt{gru}); all variants are compared under identical training/calibration protocol
(Section~\ref{sec:ablations}). Epistemic uncertainty at inference is estimated via MC-Dropout,
which in this architecture activates only the explicit post-fusion and regression-head dropout
layers -- the LSTM branches' own internal recurrent dropout is a separate mechanism tied to the
LSTM module's training-mode flag and is not re-activated during MC sampling; we state this scope
precisely rather than describing dropout as active ``throughout the network'' during MC sampling.

\subsection{Conformal calibration variants}
We compare four split-conformal variants, all under a nominal $1-\alpha=90\%$ target: \emph{pooled}
(one calibration quantile across all seven datasets), \emph{groupwise/Mondrian} (a separate
quantile per dataset, correcting for the small-group coverage failure a single pooled quantile
produces when group sizes differ by orders of magnitude), \emph{normalized groupwise} (groupwise
calibration additionally scaled per-sample by that sample's own MC-Dropout uncertainty estimate,
giving genuinely sample-varying interval widths rather than one fixed width per dataset), and the
\emph{hierarchical (empirical-Bayes)} partial-pooling calibration described below. All four require
calibration and test nonconformity scores to come from the \emph{same} predictor (MC-Dropout mean
vs.\ a deterministic forward pass) to satisfy the exchangeability assumption split-conformal
prediction relies on; we found and fixed one script (leave-one-dataset-out evaluation) that
violated this requirement (Section~\ref{sec:lodo}).

\subsection{Hierarchical (empirical-Bayes) calibration}
\label{sec:hierarchical}
Rather than treating each dataset's calibration quantile as independent (groupwise) or identical
(pooled), we estimate a DerSimonian-Laird random-effects model across the seven per-dataset
quantiles: a between-dataset variance $\hat\tau^2$ and a shrinkage weight
$w_g=\hat\tau^2/(\hat\tau^2+\hat\sigma_g^2)$ pulling each dataset's quantile toward a common mean in
proportion to how much of the observed between-dataset spread is attributable to genuine
heterogeneity versus sampling noise in that dataset's own calibration set. (A plausible alternative
formula, $w_g=n_g/(n_g+\hat\tau^2/\hat\sigma_g^2)$, does not have this interpretation and was
confirmed, during this audit, to be present in a superseded implementation that is not used for any
reported result; the correct formula above is validated independently against hand-calculated cases
and a from-scratch reimplementation, Appendix~\ref{app:verification}.) Per-dataset calibration
variance $\hat\sigma_g^2$ is estimated via a cell-clustered bootstrap (resampling whole cells, not
windows, since windows from one cell are correlated draws, not independent trials).

\subsection{Heterogeneity-guided conformal prediction (HGC-CP)}
\label{sec:hgccp}
HGC-CP assigns each calibration point a weight via a softmax over negative maximum mean discrepancy
(MMD, RBF-kernel, median-heuristic bandwidth) between the target dataset's representation and each
source dataset's representation, computes the Tibshirani et al.\ corrected weighted-conformal
quantile (normalizing over calibration points \emph{and} the test point's own implicit weight,
rather than over calibration points alone, which under-covers in the small-$n$ regime this method
targets), and applies a bounded heterogeneity-inflation term scaled by the distance to the nearest
source. We select $(\lambda,\beta)$ -- the MMD temperature and inflation strength -- via a nested
search that rotates the six non-target datasets as pseudo-targets, never touching the real target's
test set during selection (Section~\ref{sec:hgccp-results}); as the validity analysis in
\citet{barber2023beyond} makes explicit, a heuristic (non-likelihood-ratio) weight does not inherit
an exact finite-sample coverage guarantee merely by using the corrected quantile formula, so we
report HGC-CP results as an empirical, exploratory calibration method validated the same way as our
hierarchical calibration -- by measured coverage under a held-out protocol -- not as a proven
guarantee.

\subsection{Leave-one-dataset-out (model-level) generalization}
\label{sec:lodo}
We distinguish explicitly between \emph{calibration-level} generalization (the model is trained on
all seven datasets; only the calibration data for a given dataset is withheld and the calibration
quantile for that dataset is estimated, e.g.\ via partial pooling, from the other six) and
\emph{model-level} generalization (the entire dataset, including its contribution to model
training, is withheld -- leave-one-dataset-out, LODO). These measure fundamentally different
things and must not be conflated: a dataset with excellent calibration-level coverage under
hierarchical pooling can still show poor LODO coverage, because the model itself has never learned
that dataset's degradation distribution at all. We report both throughout and never describe a LODO
result as reflecting calibration-level generalization or vice versa.

\subsection{Statistical comparison protocol}
Where we compare architectures or calibration strategies across datasets with a formal test, we use
a one-sided Wilcoxon signed-rank test at the \emph{dataset} level ($n=6$ paired comparisons across
the six datasets with enough cells for both arms to be defined), not at the level of individual
architecture$\times$dataset cells, because windows -- and indeed whole cells -- are shared across
every architecture compared on the same dataset, violating the independence assumption a
per-cell-pair test would require. We report the resulting test's exact minimum attainable
one-sided $p$-value ($1/2^6=0.0156$) alongside the result, rather than only the $p$-value, so a
reader can judge how much evidential weight six paired comparisons can actually carry. Confidence
intervals for per-dataset coverage are reported both at the window level (a naive binomial/
Clopper-Pearson interval, which we found and corrected one script for silently treating windows as
independent trials; Section~\ref{sec:coverage-ci}) and via a cell-clustered bootstrap (resampling
whole cells with replacement), since the two can differ by several-fold for datasets with few
distinct test cells.

\section{Results}
\label{sec:results}

Every number in this section is produced by \texttt{make\_results\_registry.py} from the saved
model/script outputs listed in Appendix~\ref{app:verification}; none is hand-typed.

\subsection{Main model}
The full multi-scale BiLSTM-with-attention model, under groupwise conformal calibration, achieves
pooled test MAE \MainModelMAE{} (RMSE \MainModelRMSE{}) and overall coverage \MainModelCoverage\%
against a 90\% nominal target. Per-dataset coverage ranges from \CoverageNASA\% (NASA, $n=3$ test
windows -- too few to be a meaningful coverage estimate on its own) and \CoverageCALCE\% (CALCE)
up to \CoverageOxford\% (Oxford) and \CoverageXJTU\% (XJTU), with MIT (\CoverageMIT\%) and RWTH
(\CoverageRWTH\%) close to nominal.

\subsection{Architecture ablations}
\label{sec:ablations}
Table~\ref{tab:ablations} compares the full model against four ablations under identical training
and calibration protocol. Removing attention pooling (\texttt{no\_attention}) increases MAE from
\AblationFullMAE{} to \AblationNoAttentionMAE; removing the multi-scale branch
(\texttt{no\_multiscale}) increases it to \AblationNoMultiscaleMAE; the plain BiLSTM and GRU
variants increase it further still (\AblationPlainLstmMAE, \AblationGruMAE). All five variants
achieve coverage within a few points of the 90\% nominal target (\AblationFullCoverage\% to
\AblationGruCoverage\%), so the architectural choices affect point-prediction accuracy more than
calibrated coverage validity -- expected, since conformal calibration is recalibrated per variant
and corrects for whatever point-prediction error that variant produces.

\begin{table}[h]
\centering
\caption{Architecture ablations (groupwise conformal, seed 42).}
\label{tab:ablations}
\begin{tabular}{lccc}
\toprule
Variant & MAE & RMSE & Coverage (\%) \\
\midrule
Full (multi-scale + attention) & \AblationFullMAE & \AblationFullRMSE & \AblationFullCoverage \\
No attention (mean pooling)    & \AblationNoAttentionMAE & \AblationNoAttentionRMSE & \AblationNoAttentionCoverage \\
No multi-scale branch          & \AblationNoMultiscaleMAE & \AblationNoMultiscaleRMSE & \AblationNoMultiscaleCoverage \\
Plain BiLSTM                   & \AblationPlainLstmMAE & \AblationPlainLstmRMSE & \AblationPlainLstmCoverage \\
GRU                            & \AblationGruMAE & \AblationGruRMSE & \AblationGruCoverage \\
\bottomrule
\end{tabular}
\end{table}

\subsection{Baseline comparison}
Linear regression, SVR, and XGBoost baselines (Table~\ref{tab:baselines}) are fit on the same
features and splits. XGBoost (\BaselineXgboostMAE) and SVR (\BaselineSvrMAE) are competitive with
the full neural model (\AblationFullMAE) on pooled point-prediction MAE; linear regression is
clearly worse (\BaselineLinearMAE). We report conformal coverage for all three baselines, not only
the neural model, under the same predictor-consistent protocol
(Section~\ref{sec:baseline-conformal}), since a fair baseline comparison requires comparing
calibrated intervals, not only point accuracy.

\begin{table}[h]
\centering
\caption{Point-prediction baselines, pooled test set.}
\label{tab:baselines}
\begin{tabular}{lcc}
\toprule
Model & MAE & RMSE \\
\midrule
Linear regression & \BaselineLinearMAE & \BaselineLinearRMSE \\
SVR                & \BaselineSvrMAE & \BaselineSvrRMSE \\
XGBoost             & \BaselineXgboostMAE & \BaselineXgboostRMSE \\
Full model (ours)   & \AblationFullMAE & \AblationFullRMSE \\
\bottomrule
\end{tabular}
\end{table}

\subsection{Baseline conformal coverage}
\label{sec:baseline-conformal}
We extended the baseline scripts to save validation-set predictions (previously only test
predictions were saved, which made conformal calibration for these baselines impossible) and
computed pooled and groupwise conformal coverage for linear/SVR/XGBoost and, without retraining
(reusing existing checkpoints), the plain-LSTM and GRU ablations. Every architecture's groupwise
coverage lands close to the 90\% nominal target (89--91\% overall across all five architectures),
while POOLED coverage is markedly worse and more variable across architectures, underscoring that
groupwise calibration -- not the choice of point-prediction model -- is what fixes the small-group
coverage failure a single pooled quantile produces when dataset sizes differ by orders of
magnitude.

\subsection{Window-size sensitivity}
Reducing the input window from 20 to 10 cycles increases MAE from \WindowTwentyMAE{} to
\WindowTenMAE; increasing it to 30 or 50 cycles decreases MAE further, to \WindowThirtyMAE{} and
\WindowFiftyMAE{} respectively. Longer input windows give the model more history to extrapolate
from, at the cost of excluding any cell with fewer total cycles than the window length -- we do
not extend this sweep beyond 50 cycles, since NASA and Oxford's shortest cells begin to be excluded
from the ablation entirely at longer windows (visible already in the window-size results, where
NASA/Oxford per-dataset rows drop out at window=30/50; see \texttt{window\_30\_results.json} /
\texttt{window\_50\_results.json}).

\subsection{Multi-seed stability}
Across three model-initialization seeds (42, 123, 2024; cell split held fixed), MAE is
\MultiseedMAEMean{} $\pm$ \MultiseedMAEStd{} and overall coverage is \MultiseedCoverageMean\%
$\pm$ \MultiseedCoverageStd{} percentage points -- the model's point-prediction accuracy and
calibrated coverage are both stable to model-initialization randomness alone (this is separate
from, and does not substitute for, the cell-level SPLIT variability K6 addresses; see
Section~\ref{sec:k6}).

\subsection{Persistence baseline}
\label{sec:persistence}
Following the recommendation that naive forecasts be checked explicitly rather than assumed beaten
\citep{beck2025naiveforecast}, Table~\ref{tab:persistence} compares the trained model's real-unit
MAE against a trivial one-step persistence forecast (predict next SOH equals last observed SOH),
per dataset. The model does \emph{not} beat persistence on MIT or RWTH (\ModelBeatsPersistenceMIT{}
and \ModelBeatsPersistenceRWTH{} respectively) -- these are exactly the two largest, most
densely-sampled datasets, where successive SOH readings are close enough together that the naive
``no change'' forecast is already very accurate, leaving little room for a learned model to improve
on it. On every other dataset the model does beat persistence. We report this plainly rather than
omit the two datasets where it does not hold.

\begin{table}[h]
\centering
\caption{Model vs. persistence baseline, real SOH units, per dataset.}
\label{tab:persistence}
\begin{tabular}{lccc}
\toprule
Dataset & Model MAE & Persistence MAE & Model beats persistence? \\
\midrule
MIT    & \ModelMITMAEReal & \PersistenceMITMAE & \ModelBeatsPersistenceMIT \\
RWTH   & \ModelRWTHMAEReal & \PersistenceRWTHMAE & \ModelBeatsPersistenceRWTH \\
XJTU   & \ModelXJTUMAEReal & \PersistenceXJTUMAE & \ModelBeatsPersistenceXJTU \\
BIT    & \ModelBITMAEReal & \PersistenceBITMAE & \ModelBeatsPersistenceBIT \\
CALCE  & \ModelCALCEMAEReal & \PersistenceCALCEMAE & \ModelBeatsPersistenceCALCE \\
Oxford & \ModelOxfordMAEReal & \PersistenceOxfordMAE & \ModelBeatsPersistenceOxford \\
NASA   & \ModelNASAMAEReal & \PersistenceNASAMAE & \ModelBeatsPersistenceNASA \\
\bottomrule
\end{tabular}
\end{table}

\subsection{Hierarchical (empirical-Bayes) calibration}
The DerSimonian-Laird random-effects fit across five datasets with $\geq 2$ test cells (MIT, RWTH,
XJTU, BIT, CALCE -- Oxford and NASA have too few test cells for a within-dataset variance estimate)
gives $\hat\mu=$~\HierMuHat{} and $\hat\tau^2=$~\HierTauHat. The empirical-Bayes-shrunk coverage
is close to each dataset's own-quantile coverage for MIT (\HierMITCovEBShrunk\% vs.
\HierMITCovOwn\%) and RWTH (\HierRWTHCovEBShrunk\% vs. \HierRWTHCovOwn\%), but diverges sharply for
BIT (\HierBITCovEBShrunk\% vs. \HierBITCovOwn\%) -- BIT's own-quantile calibration already achieves
good coverage, and pooling toward the cross-dataset mean actively \emph{hurts} it here, visible
directly in BIT's near-zero empirical-Bayes weight ($w_{BIT}\approx0.038$, i.e. BIT's own
calibration data is trusted almost completely over the pooled prior). This is the correct, honest
behavior of a partial-pooling estimator -- it pools MORE for datasets whose own calibration
variance is high relative to the between-dataset spread, and pools LESS for datasets (like BIT)
whose own calibration set is already informative, rather than pooling by a fixed, data-independent
amount.

\subsection{HGC-CP nested hyperparameter search}
\label{sec:hgccp-results}
Table~\ref{tab:hgccp} reports, for each of the seven datasets as the real target, the $(\lambda,
\beta)$ selected by a nested search that rotates the other six as pseudo-targets without ever
touching the real target's test set, and the resulting real-target test coverage under that
selection. No dataset selects $\beta=0$ (the brief's own prior anchor setting) -- every one prefers
$\beta\geq0.5$ -- and the resulting coverage is highly inconsistent across targets: MIT, Oxford,
and NASA land at or near 100\% (uninformatively wide intervals), while BIT reaches only
\HgcCpBITCoverage\% under its own nested-selected setting. No single $(\lambda,\beta)$ generalizes
coverage uniformly across all seven sources. We report HGC-CP's results as exploratory, not
held-out-validated, per \citet{barber2023beyond}'s validity analysis for heuristic (non-likelihood-
ratio) conformal weights.

\begin{table}[h]
\centering
\caption{HGC-CP nested-search selected hyperparameters and resulting coverage, per real target.}
\label{tab:hgccp}
\begin{tabular}{lccc}
\toprule
Real target & $\lambda$ & $\beta$ & Test coverage (\%) \\
\midrule
MIT    & \HgcCpMITLam & \HgcCpMITBeta & \HgcCpMITCoverage \\
RWTH   & \HgcCpRWTHLam & \HgcCpRWTHBeta & \HgcCpRWTHCoverage \\
XJTU   & \HgcCpXJTULam & \HgcCpXJTUBeta & \HgcCpXJTUCoverage \\
BIT    & \HgcCpBITLam & \HgcCpBITBeta & \HgcCpBITCoverage \\
CALCE  & \HgcCpCALCELam & \HgcCpCALCEBeta & \HgcCpCALCECoverage \\
Oxford & \HgcCpOxfordLam & \HgcCpOxfordBeta & \HgcCpOxfordCoverage \\
NASA   & \HgcCpNASALam & \HgcCpNASABeta & \HgcCpNASACoverage \\
\bottomrule
\end{tabular}
\end{table}

\subsection{Leave-one-dataset-out (model-level) generalization}
Table~\ref{tab:lodo} reports model-level LODO: the entire held-out dataset is excluded from
training (not merely from calibration), and the predictor-consistency fix described in
Section~\ref{sec:methods} is applied uniformly. Coverage ranges from \LodoBITCoverage\% (BIT) and
\LodoXJTUCoverage\% (XJTU) up to \LodoRWTHCoverage\% (RWTH) and \LodoMITCoverage\% (MIT). Every one
of the seven held-out configurations shows higher coverage under the predictor-consistency fix than
the pre-fix (predictor-mismatched) anchor values, confirming the fix corrected a systematic
under-coverage bug rather than an isolated one. Even so, most datasets remain far below the 90\%
nominal target under genuine model-level LODO -- this is the expected, honest result of asking a
model to generalize its calibrated uncertainty to a data source it has never seen any part of
during training, a fundamentally harder task than calibration-level partial pooling
(Section~\ref{sec:hierarchical}), and we do not describe it as reflecting calibration-level
generalization.

\begin{table}[h]
\centering
\caption{Leave-one-dataset-out (model-level) results, predictor-consistent.}
\label{tab:lodo}
\begin{tabular}{lcc}
\toprule
Held-out dataset & MAE (z-units) & Coverage (\%) \\
\midrule
NASA   & \LodoNASAMAE & \LodoNASACoverage \\
CALCE  & \LodoCALCEMAE & \LodoCALCECoverage \\
Oxford & \LodoOxfordMAE & \LodoOxfordCoverage \\
MIT    & \LodoMITMAE & \LodoMITCoverage \\
BIT    & \LodoBITMAE & \LodoBITCoverage \\
XJTU   & \LodoXJTUMAE & \LodoXJTUCoverage \\
RWTH   & \LodoRWTHMAE & \LodoRWTHCoverage \\
\bottomrule
\end{tabular}
\end{table}

\subsection{Coverage confidence intervals: window-level vs. cell-clustered}
\label{sec:coverage-ci}
We found and fixed a script that reported per-dataset coverage confidence intervals treating every
test WINDOW as an independent Bernoulli trial. For datasets with many test cells (MIT, RWTH) the
window-level and cell-clustered-bootstrap intervals are close. For XJTU (8 test cells) and CALCE (2
test cells), they differ by roughly 4--5$\times$: a naive window-level interval of
approximately [92,94]\% for XJTU's coverage widens to approximately [84.5,99.3]\% once cell
clustering is accounted for, and similarly for CALCE. Oxford and NASA have too few test cells (1
each) for a cell-clustered interval to be defined at all, and we report this explicitly rather than
substituting the (misleadingly precise) window-level number in its place.

\subsection{Dataset-quality forensics: XJTU bridging sensitivity}
\label{sec:xjtu-sensitivity}
As described in Section~\ref{sec:data-quality}, \XjtuBridgingWindowCount{} of XJTU's test windows
(\XjtuBridgingFraction\%) span a removed sentinel-value block. These bridging windows show
substantially higher point-prediction error than XJTU's other test windows (RMSE roughly 5.6$\times$
higher in z-units), shifting XJTU's pooled MAE upward by roughly 8\% relative to what it would be
if these windows were excluded. Coverage is not degraded by their presence (both subsets remain at
or above the 90\% nominal target). We report both the full-XJTU and excluding-bridging-windows
numbers rather than only the former.

\subsection{Dataset-level statistical comparison (Wilcoxon)}
Comparing four architectures (full model, linear, SVR, XGBoost) pairwise at the dataset level
($n=6$ datasets with both groupwise and pooled coverage defined), groupwise calibration is closer
to the 90\% nominal target than pooled calibration in 6 of 6 datasets (one-sided Wilcoxon
signed-rank, exact $p=0.0156$, the minimum attainable with $n=6$). We report this exact minimum
alongside the result: six paired comparisons is not a basis for a strong quantitative claim about
effect size, only a qualitatively consistent direction of effect. The corresponding
architecture$\times$dataset-pair-level comparison ($n=24$) gives a much smaller nominal $p$-value,
but is not a valid independent-samples test here -- every architecture is evaluated on the SAME six
datasets' cells, so the 24 pairs are not independent trials, and we do not report that $p$-value as
if it were.

\subsection{Repeated cell-level splits (K6)}
\label{sec:k6}
% TODO: this subsection is BLOCKED on the repeated-split experiment (4 additional cell-level split
% seeds, 43-46, groupwise/pooled/normalized_groupwise + persistence + cell-clustered bootstrap,
% model-init seed held fixed at 42) still running as of this draft. Do not state a between-split
% coverage range here until outputs/repeated_split_summary.json exists and
% make_results_registry.py has pulled its numbers into numbers.tex -- see REPRODUCIBILITY.md.
\textbf{[PENDING: repeated-split experiment in progress; this subsection will report between-split
coverage variability across cell-level split seeds 42 (baseline) and 43--46 once complete, per
K6 of the audit plan. No claim is made here in its absence.]}

\section{Discussion and limitations}
\label{sec:discussion}

% TODO: write after Section 5 is fully complete (K6 included) -- a claim-discipline pass belongs
% here, checking every claim in this section against what Results actually established, per the
% brief's "would rather conclude a claim is wrong than defend it" standard. Draft after K6 lands;
% do not write the central heterogeneity-vs-pooling claim until the repeated-split evidence is in
% hand to support (or narrow) it.

\section{Conclusion}
\label{sec:conclusion}
% TODO: write last, after Discussion -- a conclusion should not introduce new claims.

\section*{Data and code availability}
% TODO: finalize dataset licensing/citation details (flagged in REPRODUCIBILITY.md as needing a
% human verification pass) before submission. Code: this manuscript's companion repository will be
% made available upon publication (author's own prior related repository is intentionally not
% cited here, per author's instruction, pending its own update).

\appendix
\section{Independent verification of statistical components}
\label{app:verification}
% TODO: transcribe already-completed verification results (see TEST_REPORT.md) into appendix form:
% conformal.py unit tests (6/6 pass, including Monte Carlo marginal coverage and an independent
% order-statistic reimplementation), HGC-CP property tests (13/13) and K10 equivalence proof
% (250 combinations + inf-width edge case), DerSimonian-Laird hand-calculated cases (3) and
% independent from-scratch reimplementation (30 random trials), Stage 2 data-layer determinism
% check (two independent end-to-end runs, identical SHA-256 hash).

\bibliographystyle{elsarticle-harv}
\bibliography{references}

\end{document}
